% Part: first-order-logic
% Chapter: completeness
% Section: lindenbaums-lemma

\documentclass[../../../include/open-logic-section]{subfiles}

\begin{document}

\iftag{FOL}
      {\olfileid{fol}{com}{lin}}
      {\olfileid{pl}{com}{lin}}

\olsection{Lemma Lindenbaum}

\begin{explain}
Saiki kita mbuktekake lemma kang nuduhake yen sembarang himpunan
konsisten saka !!{sentence}s kamot ing sawijine himpunan ukara kang
ora mung konsisten, nanging uga !!{complete}. Bukti kasebut lumaku
kanthi nambah siji !!{sentence} saben langkah lan njamin ing saben
langkah yen himpunane tetep konsisten. Iki ditindakake supaya kanggo
saben $!A$, $!A$ utawa $\lnot !A$ ditambahake ing sawijine tahap.
Gabungan kabeh tahap ing konstruksi kasebut banjur ngemot $!A$ utawa
negasine~$\lnot !A$, lan mula jangkep. Gabungan iku uga konsisten
amarga ing saben tahap kita mesthekake supaya ora ngenalake
inkonsistensi.
\end{explain}

\begin{lem}[Lemma Lindenbaum]
\ollabel{lem:lindenbaum} Saben himpunan konsisten~$\Gamma$ ing
basa~$\Lang{L}$ bisa diambakake dadi !!a{complete} himpunan
konsisten~$\Gamma^*$.
\end{lem}

\begin{proof}
Upamana $\Gamma$ konsisten. Upamana $!A_0$, $!A_1$, \dots{} minangka
enumerasi kabeh !!{sentence}s saka~$\Lang L$. Tetepna $\Gamma_0 =
\Gamma$, lan
\[
\Gamma_{n+1} =
\begin{cases}
\Gamma_n \cup \{ !A_n \} & \textrm{yen $\Gamma_n \cup \{!A_n\}$
  konsisten;} \\
\Gamma_n \cup \{ \lnot !A_n \} & \textrm{yen ora.}
\end{cases}
\]
Tetepna $\Gamma^* = \bigcup_{n \geq 0} \Gamma_n$.

Saben $\Gamma_n$ konsisten: $\Gamma_0$ konsisten miturut definisi.
Yen $\Gamma_{n+1} = \Gamma_n \cup \{!A_n\}$, iki amarga himpunan
pungkasan kasebut konsisten. Yen ora, $\Gamma_{n+1} = \Gamma_n \cup
\{\lnot !A_n\}$. Kita kudu mbuktekake yen $\Gamma_n \cup \{\lnot
!A_n\}$ konsisten. Upamana ora konsisten. Banjur \emph{kalorone}
$\Gamma_n \cup \{!A_n\}$ lan $\Gamma_n \cup \{\lnot !A_n\}$ ora
konsisten. Iki ateges $\Gamma_n$ mesthi ora konsisten miturut
\iftag{FOL}{%
  \tagrefs{prfAX/{fol:axd:prv:prop:provability-exhaustive},
    prfSC/{fol:seq:prv:prop:provability-exhaustive},
    prfND/{fol:ntd:prv:prop:provability-exhaustive},
    prfTab/{fol:tab:prv:prop:provability-exhaustive}}}{%
  \tagrefs{prfAX/{pl:axd:prv:prop:provability-exhaustive},
    prfSC/{pl:seq:prv:prop:provability-exhaustive},
    prfND/{pl:ntd:prv:prop:provability-exhaustive},
    prfTab/{pl:tab:prv:prop:provability-exhaustive}}},
kang cengkah karo hipotesis induksi.

Kanggo saben~$n$ lan saben $i < n$, $\Gamma_i \subseteq \Gamma_n$.
Iki bisa didudut kanthi induksi prasaja marang~$n$. Kanggo $n=0$,
ora ana $i < 0$, mula pratelan kasebut bener kanthi langsung. Kanggo
langkah induksi, upamana pratelan kasebut bener kanggo~$n$. Kita
nuduhake yen $i < n+1$ mula $\Gamma_i \subseteq \Gamma_{n+1}$. Miturut
konstruksi, kita duwe $\Gamma_{n+1} = \Gamma_n \cup \{!A_n\}$ utawa
$= \Gamma_n \cup \{\lnot !A_n\}$. Dadi $\Gamma_n \subseteq
\Gamma_{n+1}$. Yen $i < n+1$, mula $\Gamma_i \subseteq \Gamma_n$
miturut hipotesis induksi (yen $i < n$) utawa kasunyatan langsung yen
$\Gamma_n \subseteq \Gamma_n$ (yen $i = n$). Kita entuk $\Gamma_i
\subseteq \Gamma_{n+1}$ miturut transitivitas saka~$\subseteq$.

Saka iki bisa didudut yen $\Gamma^*$ konsisten. Katrangane mangkene:
upamana $\Gamma' \subseteq \Gamma^*$ winates. Cathetan koreksi sumber
OLPL-046: yen subhimpunan winates iki kosong, konsistensine langsung;
pilihan indeks paling gedhe ing ukara sabanjure mung kanggo perkara
kang ora kosong. Saben $!B \in \Gamma'$ uga ana ing~$\Gamma_i$ kanggo
sawijine~$i$. Upamana $n$ minangka sing paling gedhe ing antarane
indeks-indeks kasebut. Amarga $\Gamma_i \subseteq \Gamma_n$ yen $i
\le n$, saben $!B \in \Gamma'$ uga $\in \Gamma_n$, yaiku, $\Gamma'
\subseteq \Gamma_n$, lan $\Gamma_n$~konsisten. Dadi, saben subhimpunan
winates $\Gamma' \subseteq \Gamma^*$ konsisten. Miturut
\iftag{FOL}{%
  \tagrefs{prfAX/{fol:axd:ptn:prop:proves-compact},
    prfSC/{fol:seq:ptn:prop:proves-compact},
    prfND/{fol:ntd:ptn:prop:proves-compact},
    prfTab/{fol:tab:ptn:prop:proves-compact}}}{%
  \tagrefs{prfAX/{pl:axd:ptn:prop:proves-compact},
    prfSC/{pl:seq:ptn:prop:proves-compact},
    prfND/{pl:ntd:ptn:prop:proves-compact},
    prfTab/{pl:tab:ptn:prop:proves-compact}}}, $\Gamma^*$ iku
  konsisten.

Saben !!{sentence} saka $\Frm[L]$ muncul ing daftar kang digunakake
kanggo nemtokake~$\Gamma^*$. Yen $!A_n \notin \Gamma^*$, iku amarga
$\Gamma_n \cup \{!A_n\}$ ora konsisten. Nanging banjur $\lnot !A_n
\in \Gamma^*$, dadi $\Gamma^*$ iku !!{complete}.
\end{proof}

\end{document}
